% Genitivartikel und Genitivendungen
% Stilbasis: erklaerung_vorlage.tex und tabelle_vorlage.tex

\documentclass[a4paper,11pt]{article}

\usepackage[a4paper,margin=2cm]{geometry}
\usepackage[T1]{fontenc}
\usepackage[utf8]{inputenc}
\usepackage[ngerman,english]{babel}
\usepackage{lmodern}
\usepackage{microtype}
\usepackage[table]{xcolor}
\usepackage{parskip}
\usepackage{enumitem}
\usepackage[most]{tcolorbox}
\usepackage{titlesec}
\usepackage[hidelinks]{hyperref}
\usepackage{graphicx}
\usepackage{array}
\usepackage{longtable}
\usepackage{booktabs}
\graphicspath{{./}{/mnt/data/}}

\definecolor{HeaderBlueDark}{RGB}{70,110,170}
\definecolor{RowBlueLight}{RGB}{235,243,255}
\definecolor{RowBlueDark}{RGB}{220,233,250}
\definecolor{SoftGray}{RGB}{90,90,90}

\setlength{\parindent}{0pt}
\setlist[itemize]{leftmargin=1.4em,itemsep=0.35em,topsep=0.35em}
\linespread{1.06}
\renewcommand{\arraystretch}{1.5}
\setlength{\tabcolsep}{6pt}

\titleformat{\section}[block]
{\Large\bfseries\color{HeaderBlueDark}}
{}{0pt}{}
\titlespacing{\section}{0pt}{1.3em}{0.8em}

\titleformat{\subsection}[block]
{\large\bfseries\color{HeaderBlueDark}}
{}{0pt}{}
\titlespacing{\subsection}{0pt}{1.0em}{0.6em}

\newtcolorbox{exbox}{enhanced,breakable,colback=RowBlueLight,colframe=RowBlueDark,boxrule=0.4pt,arc=1.5mm,left=1.2mm,right=1.2mm,top=1mm,bottom=1mm,title={Beispiele \textnormal{(Examples)}},fonttitle=\bfseries,colbacktitle=RowBlueDark,coltitle=black}

\NewDocumentEnvironment{grammarentry}{mm+b}{%
    \begin{tcolorbox}[enhanced,breakable,colback=white,colframe=HeaderBlueDark,boxrule=0.6pt,arc=1.5mm,left=1.5mm,right=1.5mm,top=1mm,bottom=1mm,colbacktitle=HeaderBlueDark,coltitle=white,fonttitle=\bfseries,title={#1\hfill\normalfont\footnotesize #2}]
    #3
    \end{tcolorbox}
}{}

\newcommand{\GermanText}[1]{\textbf{Deutsch: }#1\par}
\newcommand{\EnglishText}[1]{{\small\color{SoftGray}\textbf{English: }#1\par}}
\newcommand{\ExampleItem}[2]{\item \textbf{#1}\\{\small\color{SoftGray}#2}}

\title{Genitivartikel und Genitivendungen}
\date{}

\begin{document}
    \maketitle
    \tableofcontents
    \newpage

    \section{Ausgangsbild der Lektion}

        \begin{grammarentry}{Lektion 2: Genitivartikel}{Inhalt des hochgeladenen Bildes}
            \GermanText{Das folgende Bild ist die Grundlage dieser Erklärung. Alle Regeln und Beispiele aus dem Bild werden in den nächsten Abschnitten vollständig erklärt.}
            \EnglishText{The following image is the basis of this explanation. All rules and examples from the image are explained completely in the following sections.}
        \end{grammarentry}

        \begin{center}
            \includegraphics[width=0.96\textwidth]{0b0a2244-02ff-48a8-b2c2-a6467a7e510c.png}
        \end{center}

    \section{Grundbedeutung des Genitivs}

        \begin{grammarentry}{Der Genitiv}{Besitz, Zugehörigkeit und Beziehung}
            \GermanText{Der Genitiv ist der zweite Kasus. Er drückt häufig Besitz, Zugehörigkeit oder eine Beziehung zwischen zwei Nomen aus. Er beantwortet die Frage \glqq Wessen?\grqq.}
            \EnglishText{The genitive is the second grammatical case. It often expresses possession, belonging, or a relationship between two nouns. It answers the question ``Whose?''}

            \GermanText{Das Nomen, das die Person oder Sache mit Besitz oder Zugehörigkeit bezeichnet, steht im Genitiv.}
            \EnglishText{The noun that identifies the person or thing associated with possession or belonging is placed in the genitive.}
        \end{grammarentry}

        \begin{exbox}
            \begin{itemize}
                \ExampleItem{Wessen Koffer ist das? -- Das ist der Koffer des Gastes.}{Whose suitcase is that? -- That is the guest's suitcase.}
                \ExampleItem{Wessen Taschen sind das? -- Das sind die Taschen der Studentin.}{Whose bags are those? -- Those are the female student's bags.}
            \end{itemize}
        \end{exbox}

    \section{Satzstruktur}

        \begin{grammarentry}{Nomen und Genitivattribut}{Grundmuster}
            \GermanText{Das Genitivattribut steht normalerweise direkt hinter dem Nomen, das näher erklärt wird.}
            \EnglishText{The genitive attribute normally appears directly after the noun that it describes more precisely.}

            \GermanText{Grundmuster: \textbf{Nomen + Genitivartikel + Genitivnomen}.}
            \EnglishText{Basic pattern: \textbf{noun + genitive article + genitive noun}.}
        \end{grammarentry}

        \begin{exbox}
            \begin{itemize}
                \ExampleItem{Der Gast hat einen Koffer. Das ist der Koffer des Gastes.}{The guest has a suitcase. That is the guest's suitcase.}
                \ExampleItem{Die Studentin hat zwei Taschen. Das sind die Taschen der Studentin.}{The female student has two bags. Those are the female student's bags.}
                \ExampleItem{Das Mädchen hat ein Fahrrad. Das ist das Fahrrad des Mädchens.}{The girl has a bicycle. That is the girl's bicycle.}
                \ExampleItem{Die Touristin hat Gepäck. Das ist das Gepäck der Touristin.}{The female tourist has luggage. That is the female tourist's luggage.}
            \end{itemize}
        \end{exbox}

    \section{Bestimmte und unbestimmte Genitivartikel}

        \begin{grammarentry}{Artikelübersicht}{die Tabelle aus dem Bild}
            \GermanText{Maskuline und neutrale Nomen haben im Genitiv die Artikel \textbf{des} oder \textbf{eines}. Feminine Nomen und Nomen im Plural haben die Artikel \textbf{der} oder, nur im Singular, \textbf{einer}.}
            \EnglishText{Masculine and neuter nouns use \textbf{des} or \textbf{eines} in the genitive. Feminine nouns and plural nouns use \textbf{der}, or \textbf{einer} only in the feminine singular.}
        \end{grammarentry}

        \rowcolors{2}{RowBlueLight}{RowBlueDark}
        \begin{longtable}{>{\bfseries}p{3.2cm}|>{\centering\arraybackslash}p{5.1cm}|>{\centering\arraybackslash}p{5.1cm}}
            \rowcolor{HeaderBlueDark}
            \color{white}\textbf{Genus / Numerus} &
            \color{white}\textbf{bestimmt} &
            \color{white}\textbf{unbestimmt} \\
            Maskulin & des + -(e)s & eines + -(e)s \\
            Feminin & der & einer \\
            Neutrum & des + -(e)s & eines + -(e)s \\
            Plural & der & kein unbestimmter Artikel \\
        \end{longtable}

        \begin{grammarentry}{English table explanation}{Definite and indefinite forms}
            \GermanText{\textbf{bestimmt} bedeutet definite article; \textbf{unbestimmt} bedeutet indefinite article. Im Plural gibt es keinen unbestimmten Artikel.}
            \EnglishText{\textbf{Bestimmt} means definite article; \textbf{unbestimmt} means indefinite article. There is no indefinite article in the plural.}
        \end{grammarentry}

        \begin{exbox}
            \begin{itemize}
                \ExampleItem{des Mannes / eines Mannes}{of the man / of a man}
                \ExampleItem{der Frau / einer Frau}{of the woman / of a woman}
                \ExampleItem{des Kindes / eines Kindes}{of the child / of a child}
                \ExampleItem{der Kinder}{of the children}
            \end{itemize}
        \end{exbox}

    \section{Die wichtigste Endungsregel}

        \begin{grammarentry}{Achtung}{nur Maskulin und Neutrum}
            \GermanText{Nur maskuline und neutrale Nomen können im Genitiv Singular die Endung \textbf{-s} oder \textbf{-es} bekommen.}
            \EnglishText{Only masculine and neuter nouns can take the ending \textbf{-s} or \textbf{-es} in the genitive singular.}

            \GermanText{Feminine Nomen und Pluralnomen bekommen normalerweise keine zusätzliche Genitivendung.}
            \EnglishText{Feminine nouns and plural nouns normally do not receive an additional genitive ending.}
        \end{grammarentry}

        \begin{exbox}
            \begin{itemize}
                \ExampleItem{der Mann -- des Mannes}{the man -- of the man}
                \ExampleItem{das Kind -- des Kindes}{the child -- of the child}
                \ExampleItem{die Frau -- der Frau}{the woman -- of the woman}
                \ExampleItem{die Kinder -- der Kinder}{the children -- of the children}
            \end{itemize}
        \end{exbox}

    \section{Wann benutzt man -s?}

        \begin{grammarentry}{Die Endung -s}{viele mehrsilbige Nomen}
            \GermanText{Viele mehrsilbige maskuline und neutrale Nomen bekommen im Genitiv nur \textbf{-s}. Das gilt besonders häufig für Nomen auf \textbf{-en}, \textbf{-el}, \textbf{-er}, \textbf{-or} und \textbf{-ling}.}
            \EnglishText{Many multisyllabic masculine and neuter nouns take only \textbf{-s} in the genitive. This is especially common with nouns ending in \textbf{-en}, \textbf{-el}, \textbf{-er}, \textbf{-or}, and \textbf{-ling}.}
        \end{grammarentry}

        \begin{exbox}
            \begin{itemize}
                \ExampleItem{des Wagens}{of the car}
                \ExampleItem{des Onkels}{of the uncle}
                \ExampleItem{des Reporters}{of the reporter}
                \ExampleItem{des Doktors}{of the doctor}
                \ExampleItem{des Frühlings}{of the spring}
                \ExampleItem{des Lebens}{of life}
                \ExampleItem{des Fensters}{of the window}
            \end{itemize}
        \end{exbox}

    \section{Wann benutzt man -es?}

        \begin{grammarentry}{Die Endung -es}{einsilbige Nomen und bestimmte Endlaute}
            \GermanText{Viele einsilbige maskuline und neutrale Nomen bekommen im Genitiv \textbf{-es}.}
            \EnglishText{Many one-syllable masculine and neuter nouns take \textbf{-es} in the genitive.}

            \GermanText{Die Endung \textbf{-es} steht außerdem gewöhnlich bei Nomen auf \textbf{-s}, \textbf{-ss}, \textbf{-ß}, \textbf{-sch}, \textbf{-z} und \textbf{-tz}.}
            \EnglishText{The ending \textbf{-es} is also normally used with nouns ending in \textbf{-s}, \textbf{-ss}, \textbf{-ß}, \textbf{-sch}, \textbf{-z}, and \textbf{-tz}.}
        \end{grammarentry}

        \begin{exbox}
            \begin{itemize}
                \ExampleItem{des Buches}{of the book}
                \ExampleItem{des Flusses}{of the river}
                \ExampleItem{des Fußes}{of the foot}
                \ExampleItem{des Schreibtisches}{of the desk}
                \ExampleItem{des Fußballplatzes}{of the football field}
                \ExampleItem{der Bus -- des Busses}{the bus -- of the bus}
            \end{itemize}
        \end{exbox}

        \begin{grammarentry}{Praktischer Hinweis}{-s und -es}
            \GermanText{Bei manchen Nomen sind sowohl \textbf{-s} als auch \textbf{-es} möglich. Die Form mit \textbf{-es} klingt häufig vollständiger oder formeller, besonders bei einsilbigen Nomen.}
            \EnglishText{With some nouns, both \textbf{-s} and \textbf{-es} are possible. The form with \textbf{-es} often sounds fuller or more formal, especially with one-syllable nouns.}

            \GermanText{Beim Lernen ist die folgende Faustregel hilfreich: mehrsilbig meistens \textbf{-s}; einsilbig meistens \textbf{-es}.}
            \EnglishText{The following rule of thumb is useful when learning: multisyllabic nouns usually take \textbf{-s}; one-syllable nouns usually take \textbf{-es}.}
        \end{grammarentry}

    \section{Genitiv mit Personennamen}

        \begin{grammarentry}{Name + s + Nomen}{kein Artikel}
            \GermanText{Bei Personennamen steht der Name normalerweise vor dem Nomen und bekommt die Endung \textbf{-s}. Vor dem Namen steht dabei kein Artikel.}
            \EnglishText{With personal names, the name normally appears before the noun and takes the ending \textbf{-s}. No article is placed before the name in this structure.}
        \end{grammarentry}

        \begin{exbox}
            \begin{itemize}
                \ExampleItem{Julia hat eine Wohnung. Das ist Julias Wohnung.}{Julia has an apartment. That is Julia's apartment.}
                \ExampleItem{Paul hat ein Fahrrad. Das ist Pauls Fahrrad.}{Paul has a bicycle. That is Paul's bicycle.}
                \ExampleItem{Goethes Werke sind weltberühmt.}{Goethe's works are world-famous.}
            \end{itemize}
        \end{exbox}

        \begin{grammarentry}{Apostroph bei Namen}{wichtiger Unterschied zum Englischen}
            \GermanText{Vor dem Genitiv-s steht im Deutschen normalerweise \textbf{kein Apostroph}: \textbf{Julias Wohnung}, nicht \textbf{Julia's Wohnung}.}
            \EnglishText{German normally uses \textbf{no apostrophe} before the genitive \textbf{-s}: \textbf{Julias Wohnung}, not \textbf{Julia's Wohnung}.}

            \GermanText{Endet ein Name bereits auf einen s-Laut, wird nur ein Apostroph gesetzt: \textbf{Max' Auto}, \textbf{Hans' Bruder}.}
            \EnglishText{If a name already ends in an s-sound, only an apostrophe is added: \textbf{Max' Auto}, \textbf{Hans' Bruder}.}
        \end{grammarentry}

    \section{Kein-Artikel und Possessivartikel}

        \begin{grammarentry}{Endungen der Ein-Wörter}{kein, mein, dein, sein, ihr, unser, euer, Ihr}
            \GermanText{Der Negativartikel \textbf{kein} und die Possessivartikel haben im Genitiv dieselben Endungen wie der unbestimmte Artikel: Maskulin und Neutrum \textbf{-es}, Feminin und Plural \textbf{-er}.}
            \EnglishText{The negative article \textbf{kein} and possessive articles use the same genitive endings as the indefinite article: masculine and neuter \textbf{-es}, feminine and plural \textbf{-er}.}
        \end{grammarentry}

        \rowcolors{2}{RowBlueLight}{RowBlueDark}
        \begin{longtable}{>{\bfseries}p{3.2cm}|>{\centering\arraybackslash}p{3.2cm}|>{\centering\arraybackslash}p{3.2cm}|>{\centering\arraybackslash}p{3.2cm}}
            \rowcolor{HeaderBlueDark}
            \color{white}\textbf{Genus / Numerus} &
            \color{white}\textbf{kein} &
            \color{white}\textbf{mein} &
            \color{white}\textbf{dein} \\
            Maskulin & keines & meines & deines \\
            Feminin & keiner & meiner & deiner \\
            Neutrum & keines & meines & deines \\
            Plural & keiner & meiner & deiner \\
        \end{longtable}

        \begin{exbox}
            \begin{itemize}
                \ExampleItem{die Meinung meines Bruders}{my brother's opinion}
                \ExampleItem{die Farbe deiner Jacke}{the color of your jacket}
                \ExampleItem{das Ende keines Films}{the ending of no film}
                \ExampleItem{die Namen unserer Gäste}{the names of our guests}
            \end{itemize}
        \end{exbox}

    \section{Adjektive nach Genitivartikeln}

        \begin{grammarentry}{Artikel + Adjektiv + Nomen}{Adjektivendung -en}
            \GermanText{Nach einem bestimmten Artikel, einem unbestimmten Artikel, einem Negativartikel oder einem Possessivartikel bekommt das Adjektiv im Genitiv normalerweise die Endung \textbf{-en}.}
            \EnglishText{After a definite article, indefinite article, negative article, or possessive article, an adjective normally takes the ending \textbf{-en} in the genitive.}
        \end{grammarentry}

        \begin{exbox}
            \begin{itemize}
                \ExampleItem{das Auto des alten Mannes}{the old man's car}
                \ExampleItem{die Tür eines neuen Hauses}{the door of a new house}
                \ExampleItem{der Name einer bekannten Autorin}{the name of a well-known female author}
                \ExampleItem{die Stimmen unserer jungen Mitglieder}{the voices of our young members}
            \end{itemize}
        \end{exbox}

    \section{Wichtige Sonderfälle bei Nomen}

        \begin{grammarentry}{N-Deklination}{schwache maskuline Nomen}
            \GermanText{Viele schwache maskuline Nomen bekommen im Genitiv nicht \textbf{-s} oder \textbf{-es}, sondern die Endung \textbf{-n} oder \textbf{-en}.}
            \EnglishText{Many weak masculine nouns do not take \textbf{-s} or \textbf{-es} in the genitive; instead, they take \textbf{-n} or \textbf{-en}.}
        \end{grammarentry}

        \begin{exbox}
            \begin{itemize}
                \ExampleItem{der Student -- des Studenten}{the student -- of the student}
                \ExampleItem{der Mensch -- des Menschen}{the person -- of the person}
                \ExampleItem{der Junge -- des Jungen}{the boy -- of the boy}
                \ExampleItem{der Präsident -- des Präsidenten}{the president -- of the president}
            \end{itemize}
        \end{exbox}

        \begin{grammarentry}{Besondere Formen}{-ns und andere feste Formen}
            \GermanText{Einige Nomen haben besondere Genitivformen, die man als Wortform lernen sollte.}
            \EnglishText{Some nouns have special genitive forms that should be learned as individual word forms.}
        \end{grammarentry}

        \begin{exbox}
            \begin{itemize}
                \ExampleItem{der Name -- des Namens}{the name -- of the name}
                \ExampleItem{der Gedanke -- des Gedankens}{the thought -- of the thought}
                \ExampleItem{der Wille -- des Willens}{the will -- of the will}
                \ExampleItem{das Herz -- des Herzens}{the heart -- of the heart}
            \end{itemize}
        \end{exbox}

    \section{Genitiv nach Präpositionen}

        \begin{grammarentry}{Häufige Genitivpräpositionen}{wegen, trotz, während und andere}
            \GermanText{Nach mehreren Präpositionen steht in der Standardsprache der Genitiv. Häufige Beispiele sind \textbf{wegen}, \textbf{trotz}, \textbf{während}, \textbf{statt / anstatt}, \textbf{innerhalb}, \textbf{außerhalb}, \textbf{aufgrund} und \textbf{bezüglich}.}
            \EnglishText{In standard German, several prepositions require the genitive. Common examples are \textbf{wegen}, \textbf{trotz}, \textbf{während}, \textbf{statt / anstatt}, \textbf{innerhalb}, \textbf{außerhalb}, \textbf{aufgrund}, and \textbf{bezüglich}.}
        \end{grammarentry}

        \begin{exbox}
            \begin{itemize}
                \ExampleItem{Wegen des Regens bleiben wir zu Hause.}{Because of the rain, we are staying at home.}
                \ExampleItem{Trotz der Schmerzen arbeitet er weiter.}{Despite the pain, he continues working.}
                \ExampleItem{Während des Unterrichts darf man nicht telefonieren.}{During the lesson, one must not use the phone.}
                \ExampleItem{Außerhalb der Stadt sind die Wohnungen günstiger.}{Outside the city, apartments are less expensive.}
            \end{itemize}
        \end{exbox}

    \section{Genitiv oder von + Dativ?}

        \begin{grammarentry}{Standardsprache und Umgangssprache}{Bedeutung bleibt ähnlich}
            \GermanText{In der gesprochenen Sprache wird ein Genitivattribut häufig durch \textbf{von + Dativ} ersetzt. In formellen Texten und in Prüfungen ist der Genitiv jedoch besonders wichtig.}
            \EnglishText{In spoken German, a genitive attribute is often replaced by \textbf{von + dative}. In formal texts and examinations, however, the genitive is especially important.}
        \end{grammarentry}

        \begin{exbox}
            \begin{itemize}
                \ExampleItem{das Auto des Mannes}{the man's car; standard genitive}
                \ExampleItem{das Auto von dem Mann}{the car belonging to the man; common in spoken German}
                \ExampleItem{die Meinung meiner Schwester}{my sister's opinion; standard genitive}
                \ExampleItem{die Meinung von meiner Schwester}{the opinion of my sister; common in spoken German}
            \end{itemize}
        \end{exbox}

        \begin{grammarentry}{Wichtige Einschränkung}{nicht immer austauschbar}
            \GermanText{Nach Präpositionen, die den Genitiv verlangen, sollte man in der Standardsprache den Genitiv verwenden: \textbf{wegen des Wetters}, \textbf{trotz der Probleme}.}
            \EnglishText{After prepositions that require the genitive, the genitive should be used in standard German: \textbf{wegen des Wetters}, \textbf{trotz der Probleme}.}
        \end{grammarentry}

    \section{Typische Fehler}

        \begin{grammarentry}{Fehler vermeiden}{Artikel, Endung und Apostroph}
            \GermanText{Bei maskulinen und neutralen Nomen muss man sowohl den Artikel als auch meistens das Nomen verändern.}
            \EnglishText{With masculine and neuter nouns, both the article and usually the noun itself must be changed.}
        \end{grammarentry}

        \begin{exbox}
            \begin{itemize}
                \ExampleItem{Falsch: der Koffer des Gast -- Richtig: der Koffer des Gastes}{Wrong: of the guest without a noun ending -- Correct: des Gastes}
                \ExampleItem{Falsch: die Farbe der Haus -- Richtig: die Farbe des Hauses}{Wrong article -- Correct: des Hauses}
                \ExampleItem{Falsch: das Fahrrad des Frau -- Richtig: das Fahrrad der Frau}{Wrong article -- Correct: der Frau}
                \ExampleItem{Falsch: Julia's Wohnung -- Richtig: Julias Wohnung}{Wrong German apostrophe -- Correct without an apostrophe}
                \ExampleItem{Falsch in der Standardsprache: wegen dem Regen -- Richtig: wegen des Regens}{Non-standard case after wegen -- Standard: genitive}
            \end{itemize}
        \end{exbox}

    \section{Entscheidungsschritte}

        \begin{grammarentry}{So bildest du den Genitiv}{Schritt für Schritt}
            \GermanText{1. Bestimme Genus und Numerus des Genitivnomens.}
            \EnglishText{1. Determine the gender and number of the noun in the genitive.}

            \GermanText{2. Wähle den richtigen Artikel: \textbf{des / eines} für Maskulin und Neutrum, \textbf{der / einer} für Feminin und \textbf{der} für Plural.}
            \EnglishText{2. Choose the correct article: \textbf{des / eines} for masculine and neuter, \textbf{der / einer} for feminine, and \textbf{der} for plural.}

            \GermanText{3. Ergänze bei einem normalen maskulinen oder neutralen Nomen im Singular \textbf{-s} oder \textbf{-es}.}
            \EnglishText{3. Add \textbf{-s} or \textbf{-es} to a regular masculine or neuter singular noun.}

            \GermanText{4. Prüfe, ob das Nomen ein schwaches maskulines Nomen oder eine besondere Form ist.}
            \EnglishText{4. Check whether the noun is a weak masculine noun or has a special form.}

            \GermanText{5. Ein Adjektiv nach einem Genitivartikel bekommt normalerweise \textbf{-en}.}
            \EnglishText{5. An adjective after a genitive article normally takes \textbf{-en}.}
        \end{grammarentry}

    \section{Kurze Gesamtübersicht}

        \rowcolors{2}{RowBlueLight}{RowBlueDark}
        \begin{longtable}{>{\bfseries}p{3.5cm}|p{5.2cm}|p{5.2cm}}
            \rowcolor{HeaderBlueDark}
            \color{white}\textbf{Kategorie} &
            \color{white}\textbf{Deutsche Regel} &
            \color{white}\textbf{English rule} \\
            Maskulin & des / eines + Nomen mit -(e)s & des / eines + noun with -(e)s \\
            Feminin & der / einer; Nomen unverändert & der / einer; noun unchanged \\
            Neutrum & des / eines + Nomen mit -(e)s & des / eines + noun with -(e)s \\
            Plural & der; Nomen normalerweise unverändert & der; noun normally unchanged \\
            Personennamen & Name + s vor dem Nomen & name + s before the noun \\
            Adjektiv & nach dem Artikel normalerweise -en & normally -en after the article \\
            N-Deklination & -n oder -en statt -(e)s & -n or -en instead of -(e)s \\
        \end{longtable}

    \section{Zusammenfassung}

        \begin{grammarentry}{Merksatz}{kurz und wichtig}
            \GermanText{\textbf{des / eines} stehen bei maskulinen und neutralen Nomen; das Nomen bekommt meistens \textbf{-s} oder \textbf{-es}. \textbf{der / einer} stehen bei femininen Nomen; \textbf{der} steht im Plural.}
            \EnglishText{\textbf{Des / eines} are used with masculine and neuter nouns; the noun usually takes \textbf{-s} or \textbf{-es}. \textbf{Der / einer} are used with feminine nouns; \textbf{der} is used in the plural.}

            \GermanText{Mehrsilbige Nomen bekommen häufig \textbf{-s}; einsilbige Nomen und Nomen mit bestimmten Endlauten bekommen häufig \textbf{-es}. Schwache maskuline Nomen bilden eine wichtige Ausnahme.}
            \EnglishText{Multisyllabic nouns often take \textbf{-s}; one-syllable nouns and nouns with certain final sounds often take \textbf{-es}. Weak masculine nouns are an important exception.}

            \GermanText{Bei Personennamen gilt normalerweise: \textbf{Name + s + Nomen}, ohne Apostroph.}
            \EnglishText{With personal names, the normal pattern is: \textbf{name + s + noun}, without an apostrophe.}
        \end{grammarentry}

\end{document}
